%%%%%%%%%%%%%%%%%%%%%%%%%%%%%%%%%%%%%%%%%%%%%%%%%%%%%%%%%%%%
%% thesis.tex — Carbon Lock-in or Circular Transition?
%% Heterogeneity in Japan's Waste Incineration Fleet
%% and Net-Zero Compatibility
%% Compilable with: pdflatex
%%%%%%%%%%%%%%%%%%%%%%%%%%%%%%%%%%%%%%%%%%%%%%%%%%%%%%%%%%%%

\documentclass[12pt,a4paper]{article}

%% ---------- packages ----------
\usepackage[utf8]{inputenc}
\usepackage[margin=1in]{geometry}
\usepackage{amsmath}
\usepackage{graphicx}
\graphicspath{{figures/}}
\usepackage{booktabs}
\usepackage[hidelinks]{hyperref}
\usepackage[round]{natbib}
\usepackage{setspace}
\usepackage{float}
\usepackage{caption}
\usepackage{subcaption}
\usepackage{tabularx}

\onehalfspacing

%% ---------- title information ----------
\title{%
  Carbon Lock-in or Circular Transition?\\
  Heterogeneity in Japan's Waste Incineration Fleet\\
  and Net-Zero Compatibility%
}
\author{Pann Phetra \\ Student ID: 13223029 \\ Supervisor: Prof. Han Ji \\ Ritsumeikan Asia Pacific University}
\date{2026}

%% ==========================================================
\begin{document}
\maketitle

%% ==========================================================
%% ABSTRACT
%% ==========================================================
\begin{abstract}
Japan's waste incinerators are not all the same --- and that heterogeneity determines whether incineration helps or hurts the country's carbon goals. Using a 20-year facility-level panel of 23,599 observations across 2,949 incinerators (Ministry of the Environment, FY2005--2024), this thesis identifies the characteristics that predict energy recovery efficiency among Japan's power-generating facilities.

Three findings stand out. First, facility age exerts a consistent and substantial negative effect on efficiency: coefficients range from $-0.028$ to $-0.043$ per year across eight robustness specifications ($p < 0.001$ in all), and mean efficiency declines from 0.397 MWh/t at facilities aged 0--10 years to 0.179 MWh/t at facilities over 30 years old. Second, scale and capacity utilization are strong positive predictors, with each additional 100 t/day of design capacity associated with a 0.08--0.10 log-unit gain in efficiency. Third --- and most consequentially --- within-facility efficiency is essentially fixed: the ratio of within-to-total variation is approximately 0.11 (stable across subsamples), indicating that roughly 89\% of efficiency variation is between facilities and that efficiency is design-determined at construction and cannot be meaningfully improved through operational adjustment alone.

At the fleet level, consolidation has reduced facility count from 1,318 to 1,014 ($-23\%$) between 2005 and 2024, while the share of power-generating facilities has nearly doubled from 21.6\% to 41.1\%. Yet 59\% of active facilities still generate no electricity. Gross avoided CO$_2$ emissions from waste-to-energy reached approximately 4.6 million t-CO$_2$ in FY2024 (an upper bound excluding process emissions from combustion), but this aggregate masks profound heterogeneity: a small number of large, modern facilities account for most of the benefit.

These results reframe the policy challenge. Because efficiency is locked in at the design stage, fleet-wide improvement requires retirement of old facilities and targeted replacement --- not operational optimization of the existing stock.
\end{abstract}

\noindent\textbf{Keywords:} waste-to-energy, energy recovery efficiency, carbon lock-in, industrial ecology, Japan, fleet heterogeneity, infrastructure transition

\newpage
\tableofcontents
\newpage

%% ==========================================================
%% NOTE: Chapter content will be converted from Markdown drafts
%% in thesis/*.md files. This LaTeX skeleton provides the
%% structure and formatting for Overleaf compilation.
%%
%% Current status: Markdown chapters complete (7 files).
%% LaTeX conversion pending after citation verification.
%% ==========================================================

\section{Introduction}
\textit{Content pending conversion from 01-introduction.md}

\section{Literature Review}
\textit{Content pending conversion from 02-literature-review.md}

\section{Methodology}
\textit{Content pending conversion from 03-methodology.md}

\section{Results}
\textit{Content pending conversion from 04-results.md}

\section{Discussion}
\textit{Content pending conversion from 05-discussion.md}

\section{Conclusion}
\textit{Content pending conversion from 06-conclusion.md}

%% ==========================================================
%% BIBLIOGRAPHY
%% ==========================================================
%% Bibliography will be added after citation verification.
%% Use: \bibliography{references}
%% \bibliographystyle{apalike}

\end{document}
